\documentclass[10pt,a4paper]{amsart}
\usepackage[margin=19mm]{geometry}
\usepackage{amsmath,amssymb,mathtools}
\usepackage[scr=rsfs,cal=euler]{mathalfa}
\usepackage{xfrac}
\usepackage[unicode,hidelinks,bookmarks=false]{hyperref}
\newcommand{\ie}{i.e.\ }
\newcommand{\cf}{cf.\ }
\newcommand{\resp}{respectively\ }
\newcommand{\Subsection}{\subsection}
\providecommand{\nobreakdash}{\nobreak\hbox{-}}
\setlength{\emergencystretch}{2em}
\multlinegap0pt
\usepackage{bm}
\usepackage{bbm}
\usepackage{dsfont}
\usepackage{xspace}
\usepackage{cancel}
\usepackage{mathtools}
\usepackage{amsthm}
\makeatletter
\@ifundefined{theorem}{\newtheorem{theorem}{Theorem}}{}
\@ifundefined{lemma}{\newtheorem{lemma}[theorem]{Lemma}}{}
\makeatother
\title[On expansive operators that are quasisimilar to the unilater]{On expansive operators that are quasisimilar to the unilateral shift of finite multiplicity}
\author{Maria F. Gamal'}
\date{Source: arXiv:2303.16300v2 (CC BY 4.0)}
\begin{document}
\maketitle

\noindent\emph{Source and licence.} On expansive operators that are quasisimilar to the unilateral shift of finite multiplicity. Version arXiv:2303.16300v2 (CC BY 4.0), distributed under the Creative Commons Attribution 4.0 International licence, as recorded in the arXiv licence metadata for this exact version and shown on the abstract page https://arxiv.org/abs/2303.16300v2. Full rights notice: licenses/arxiv-2303.16300-CC-BY-4.0.txt.\par\smallskip

\noindent\textit{Editorial scope.} Counted unit: the Theorem whose source label is \texttt{thmmainiso} in the article (source0.tex lines 815--820, source label \texttt{thmmainiso}), delivered together with the article's own complete original proof of it (source0.tex lines 822--886, one \texttt{proof} environment of 65 lines). No mathematics is written, repaired or re-derived: statement and proof are the article's own text line for line. Required context retained uncounted: lem11 (lines 727--731); lem12 (lines 746--752); lem13 (lines 769--775). Every definition, display and label of those lines is retained unchanged. Omitted from this package and not redistributed: 1--726, 732--745, 753--768, 776--814, 821--821, 887--1876. Nothing inside a retained excerpt is omitted or altered: every retained body is delivered exactly as the article's lines are written, so no omission receipt of any kind accompanies this package. Citations: the retained excerpts carry no citation command at all, so no bibliography entry of the article is transcribed and no reference is redistributed. Reference adaptation: none was needed and none was made; every reference inside the delivered unit resolves inside it, so no unresolved reference or citation is produced. Printed slips kept literal: no printed word, symbol, hypothesis or number of the counted statement or of its proof was altered. Layout and class: the article's own class is \texttt{amsart} with packages including hyperref, amsmath, amsfonts, amssymb, amsthm, babel, enumerate, graphicx, color, amsrefs; the wrapper is the lane's pinned \texttt{amsart} editorial wrapper at 10pt on A4, which restates the theorem-like environments the retained text needs and adds the article's own macro definitions that the retained text needs (none), with extra wrapper packages (bm, bbm, dsfont, xspace, cancel, mathtools). The delivered numbering is the unit's own. Assets: no retained span contains a figure environment, a graphics-inclusion command, a PDF-page inclusion, a table or a TikZ picture; no image file is copied into this package, the package declares no asset dependency and no image file is redistributed with it. Licence: version arXiv:2303.16300v2 of this article is distributed under the Creative Commons Attribution 4.0 International licence, as recorded in the arXiv licence metadata for this exact version and shown on the abstract page https://arxiv.org/abs/2303.16300v2. A full rights notice is delivered at licenses/arxiv-2303.16300-CC-BY-4.0.txt. Characters: every retained excerpt is pure ASCII, so the delivered engine, XeLaTeX with its default Latin Modern text font, reports no missing character.
\begin{lemma} \label{lem11}Suppose that $\sigma\subset\mathbb T$, $X\in\mathcal I(U_\sigma^{-1}, S^*)$, and there exists 
 $f_1\in L^2(\sigma,m)$   such that $Xf_1=\mathbf{1}$.
 Then $\sigma=\mathbb T$ 
and \begin{equation*} U_{\mathbb T}|_{\vee_{n=0}^\infty U_{\mathbb T}^n f_1}\cong S.\end{equation*}
\end{lemma}

\begin{lemma} \label{lem12} Suppose that $N\in\mathbb N$,
  $V_+\in\mathcal L(\mathcal K_+)$ is an isometry, $\dim\ker V_+^*<\infty$, 
 $X_+\in\mathcal L(\mathcal K_+, H_N^2)$, and $S_N^*X_+V_+=X_+$.
 Let $V\in\mathcal L(\mathcal K)$ be the minimal unitary extension of $V_+$.
Then there exists $X\in\mathcal L(\mathcal K, H_N^2)$ such  $S_N^*XV=X$ and 
$X|_{\mathcal K_+}=X_+$.
\end{lemma}

\begin{lemma}\label{lem13} Let $N\in\mathbb N$. Write $L^2_{N+1}=H^2_N\oplus (H^2_-)_N\oplus L^2$. 
Let $h_0\in H^2_N$, and let $f\in L^2$ be such that $\int_{\mathbb T}\log|f|\mathrm{d}m>-\infty$. 
Set \begin{equation*}
\mathcal M=H^2_N\vee\vee_{n=0}^\infty U_{\mathbb T, N+1}^n(\overline\chi\overline h_0\oplus f).
\end{equation*}
Then  $U_{\mathbb T, N+1}|_{\mathcal M}\cong S_{N+1}$.
\end{lemma}

\begin{theorem} \label {thmmainiso}Suppose than $N\in\mathbb N$, $V_+\in\mathcal L(\mathcal K_+)$ is an a.c. isometry, 
$\mu_{V_+}\leq N$, 
$X_+\in\mathcal L(\mathcal K_+, H^2_N)$, and $S_N^*X_+V_+=X_+$. Suppose that there exist
 $\{f_j\}_{j=1}^N\subset\mathcal K_+$
such that $X_+f_j=e_j$ ($j=1,\ldots, N$). Then \begin{equation*}V_+|_{\vee_{j=1}^N\vee_{n=0}^\infty V_+^n f_j}\cong S_N.\end{equation*}
\end{theorem}

\begin{proof} The theorem will be proved using induction. 
Let $N=1$. Since there exists $f_1\in\mathcal K_+$ such that $Xf_1=e_1=\mathbf{1}$, we have 
 $\mathcal K_+\neq\{0\}$, and $V_+\cong S$ or $V_+\cong U_\sigma$ for some $\sigma\subset\mathbb T$. If $V_+\cong S$, 
the conclusion of the theorem
is fulfilled for every $0\not\equiv f_1\in\mathcal K_+$. If $V_+\cong U_\sigma$, Lemma \ref{lem11} is applied.   
Thus, if $N=1$, then the theorem is proved.

If $N\geq 1$, assume that the theorem is proved for all $1\leq k\leq N$. 
We will to prove the theorem for $N+1$. Let $X$ and $V$ be from Lemma \ref{lem12} applied to $V_+$ and $X_+$. 
Then $V$ is unitary, and  $S_{N+1}^*XV=X$. Set
\begin{equation} \mathcal M_k=\vee_{j=1}^k\vee_{n\in\mathbb Z}V^n f_j 
\ \ \  \text{ and } X_k=P_{H_k^2\oplus\{0\}}X|_{\mathcal M_k} \ \ \ (k=1,\ldots, N+1). 
\end{equation} 
Then $S_k^*X_kV|_{\mathcal M_k} =X_k$ and $X_kf_j=e_j$ for $j=1, \ldots, k$ ($k=1,\ldots, N+1$). Thus,  
$X_k$ and $V|_{\mathcal M_k}$ satisfy the assumption of the theorem. By the inductive hypothesis, 
$V|_{\vee_{j=1}^k \vee_{n=0}^\infty V^n f_j}\cong S_k$ for $k=1, \ldots, N$. Consequently,
\begin{equation}\label{mmconguu} V|_{\mathcal M_k}\cong U_{\mathbb T,k}\ \ \  (k=1, \ldots, N). \end{equation}


Taking into account relations \eqref{mmconguu} and the estimate $\mu_V\leq N+1$, and using appropriate unitary equivalence, we may assume that $V=U_{\mathbb T,N}\oplus U_\sigma$ 
for some $\sigma\subset \mathbb T$, and $\mathcal M_k=L^2_k\oplus\{0\}\subset L^2_N$ ($k=1, \ldots, N$). 

Write $S_{N+1}^*$ as $(N+1)\times (N+1)$  diagonal   matrix, whose elements on the main diagonal are 
$S^*$. Write  $V$  as $(N+1)\times (N+1)$  diagonal   matrix, whose $N$ elements on the main diagonal are  $U_{\mathbb T}$
and the ending element is $U_\sigma$. Write $X$ as a $(N+1)\times (N+1)$ matrix: 
$X=[X_{jk}]_{j,k=1,\ldots, N+1}$.  
Then  $S^*X_{jk}U_{\mathbb T}=X_{jk}$  and $S^*X_{j, N+1}U_{\sigma}=X_{j, N+1}$ for all $j=1,\ldots, N+1$, $k=1,\ldots,N$. 
Therefore, there exist 
$\psi_{jk}\in L^\infty$ such that $X_{jk}f=P_+\psi_{jk} f$ for every $f\in L^2$, and $X_{j,N+1}f=P_+\psi_{j,N+1} f$ for every $f\in L^2(\sigma,m)$ and for all $j=1,\ldots, N+1$, $k=1,\ldots,N$.

Set $\Psi=[\psi_{jk}]_{j,k=1,\ldots, N+1}$. 
For $k=1,\ldots, N$ write $f_k\in L^2_k\oplus\{0\}$ as a column whose first $k$ elements are functions from $L^2$ and other are zeros functions. 
 Write $f_{N+1}$ as a column whose first $N$ elements are functions from $L^2$ and $(N+1)$th element is a 
function from $L^2(\sigma,m)$. Set $F=[f_1,\ldots,f_{N+1}]$.
 Then $F$ is a upper-triangular $(N+1)\times (N+1)$ matrix, 
whose elements are functions from $L^2$ and $L^2(\sigma,m)$. Denote the elements from the main diagonal of $F$ by $f_{0k}$ ($k=1,\ldots, N+1$).
Then $f_{0k}\in L^2$ for $k=1,\ldots, N$,  $f_{0,N+1}\in L^2(\sigma,m)$, and $\det F=\prod_{k=1}^{N+1}f_{0k}$.

Since $Xf_j=e_j$ ($j=1,\ldots, N+1$), we have $P_+\Psi F=I_{(N+1)\times (N+1)}$. Therefore,
 there exists $(N+1)\times (N+1)$ matrix $G$, whose elements are functions from $H^2$, such that 
$\Psi F=I_{(N+1)\times (N+1)}+\overline\chi\overline G$.
Set $h=\det(I_{(N+1)\times (N+1)}+\chi G)$. Then $h\in H^{\frac{2}{N+1}}$, and $h(0)=1$. Therefore,  
$\int_{\mathbb T}\log|h|\mathrm{d}m>-\infty$. 
Set $\psi=\det \Psi$. Then $\psi\in L^\infty$. We have 
\begin{equation*}\overline h=\det(\Psi F)=\det\Psi\det F=\psi\prod_{k=1}^{N+1}f_{0k}.
\end{equation*}
Therefore, 
\begin{equation*}
\int_{\mathbb T}\log|\psi|\mathrm{d}m+\sum_{k=1}^{N+1}\int_{\mathbb T}\log|f_{0k}|\mathrm{d}m=
\int_{\mathbb T}\log|h|\mathrm{d}m>-\infty.
\end{equation*}
We obtain that $\int_{\mathbb T}\log|f_{0k}|\mathrm{d}m>-\infty$ for all $k=1,\ldots, N+1$. In particular, $\sigma=\mathbb T$ 
and $V=U_{\mathbb T,N+1}$. 

By the inductive hypothesis, $V|_{\vee_{j=1}^N \vee_{n=0}^\infty V^n f_j}\cong S_N$. We may assume that 
\begin{equation*}\vee_{j=1}^N \vee_{n=0}^\infty V^n f_j=H^2_N\oplus\{0\}\oplus\{0\}\subset 
H^2_N\oplus (H^2_-)_N\oplus L^2=L^2_{N+1}. \end{equation*}  
Note that  $f_{0,N+1}=P_{\{0\}\oplus\{0\}\oplus L^2}f_{N+1}$. 
Set $\overline\chi\overline h_0=P_{\{0\}\oplus (H^2_-)_N\oplus\{0\}}f_{N+1}$, where $h_0\in H^2_N$.  
Then \begin{equation*}\vee_{j=1}^{N+1} \vee_{n=0}^\infty V^n f_j 
=H^2_N\vee\vee_{n=0}^\infty V^n(\overline\chi\overline h_0\oplus f_{0,N+1}).
\end{equation*}
By Lemma \ref{lem13}, \begin{equation*}V|_{\vee_{j=1}^{N+1} \vee_{n=0}^\infty V^n f_j}\cong S_{N+1}.\end{equation*}
Since $V$ is a unitary extension of $V_+$, the theorem is proved. 
\end{proof}

\end{document}
